\begin{textsample}{NEG Dim 7 – global\_south\_2024\_10 – Score -7.00 – global\_south\_2024\_10/115617943.txt}  \label{ex:f7_neg_249}
People light their mobile phones and hold portraits of victims during a ceremony organised by the Jewish community in \textbf{commemoration} for the victims of the Hamas attack on October 7 last year in Buenos Aires. ( AFP ) The Tel Aviv ceremony took place in Israel ' s largest \textbf{concert} venue, where event planners had handed out some 40,000 tickets for the \textbf{commemoration} Hundreds gathered at the site of the Supernova \textbf{festival} on Monday to remember the 364 people who were killed and 60 others taken hostage by Hamas exactly one year ago.
Families and relatives of those killed or abducted during the October 7 attack on Israel gathered in Tel Aviv, where images of victims were \textbf{flashed} on screens.
The ceremony was broadcast live and interspersed with pre-recorded video clips of relatives telling the tragic stories of their deaths. ' Screams for help ' In a video broadcast at the event, Yuval Sharvit Trabelsi, who survived the Nova massacre, revealed for the first time that she witnessed rape while trying to evade Hamas terrorists."We saw \textbf{murder}, kidnappings, but the @ @ @ @ @ @ @ @ @ @ quoted as saying by the Jerusalem Post."I have never heard screams for help like the ones I heard from that woman."Trabelsi went on to recount how she smeared herself with her husband ' s blood so that the terrorists would think she was dead.
In all, more than 360 people were killed at the Nova \textbf{festival}.
The event in Tel Aviv ' s Hayarkon Park served as an alternative to the official government ceremony, with participating families hoping to avoid being embroiled in divisive Israeli politics as they marked the sombre anniversary.
The theme of the event was unity, with relatives of the country ' s Jewish, Arab, and Druze communities along with a mix of civilians and soldiers speaking about those who were killed.
One year ago, Iran-backed terrorists launched a war to wipe Israel off the map.They \textbf{slaughtered} more than a thousand innocents and took 251 hostage.
Every day since, Israel has been working to bring all the hostages home while protecting its citizens and fighting a multi-front... @ @ @ @ @ @ @ @ @ @ took place in Israel ' s largest \textbf{concert} venue, where event planners had handed out some 40,000 tickets for the \textbf{commemoration}.
However, amid a flurry of missile and rocket attacks by Iran-backed groups targeting Israel, the ceremony was restricted to family members and media only.
Around the country, small groups of Israelis gathered to watch the ceremony live on large screens in parks, schools and community centres instead. ' Will continue to fight. ' The official government event aired shortly after the Tel Aviv ceremony wrapped, with Prime Minister Benjamin Netanyahu vowing to press on fighting until achieving the"sacred mission"of the war against Hamas."As long as the enemy threatens our existence and the peace of our country, we will continue to fight.
As long as our hostages are still in Gaza, we will continue to fight,"said Netanyahu in a televised address.
PM Netanyahu :"A year ago, on October 7th, we were all \textbf{hurled} into a critical battle. @ @ @ @ @ @ @ @ @ @ 6:29 am local time with a minute ' s silence at the site of the Nova \textbf{Festival}, where the largest number of people -- an estimated 370 -- were killed.
In nearby communities, which were also attacked when hundreds of Hamas militants broke through the Gaza border fence, smaller \textbf{commemoration} events were held.
Participants remembered those who were killed or taken hostage, including some who are still being held captive in Gaza a year later.
Hamas ' s attack, which resulted in the deaths of 1,206 people, most of them civilians, according to an AFP tally based on official Israeli figures, sparked a full-scale war in Gaza that continues today.
Israel ' s retaliatory military offensive in Gaza has killed at least 41,909 people, a majority of them civilians, according to figures provided by the Palestinian territory ' s health ministry.

% matched lemmas: commemoration, concert, festival, flash, hurl, murder, slaughter
\end{textsample}
